\chapter{Estudo de caso pareado e cronologia do entorno}
\label{cap:caso}

Duas usinas do mesmo subsistema, sujeitas ao mesmo despacho, perdem frações
distintas da energia disponível. As explicações locais que a diferença sugere
são testáveis na frota, e o período anterior ao início da apuração é alcançável
pelo fator de capacidade.

\section{Dois ativos com taxas de corte distintas}
\label{sec:fig3}

\begin{ficha}
  \item[Janela]      outubro de 2024 a agosto de 2026, recorte comum ao par
  \item[Amostra]     2 ativos no par; 60 conjuntos no teste de frota
  \item[Unidade]     conjunto
  \item[Fonte]       ONS e ANEEL/SIGA
  \item[Teste]       Spearman sobre 5 covariáveis, sem correção
\end{ficha}

Sol do Cerrado tem 681,3 MW, é produtor independente, reúne 17 sociedades de
propósito específico de um mesmo grupo e conecta em barramento compartilhado de
Jaíba, em 230 kV. O conjunto Três Marias 3 tem 70,0 MW, opera sob autoprodução
de sete empresas privadas distintas e tem ponto de conexão exclusivo. Os dois
ficam em Minas Gerais, no mesmo subsistema. A janela comum foi recortada
programaticamente.

\begin{figure*}[t]
\centering
\includegraphics[width=\linewidth]{../figures/fig3_estudo_caso.pdf}
\caption{\textbf{Dois ativos solares em Minas Gerais diferem 3,0 vezes na taxa
de corte, e a diferença não se reproduz como padrão no parque.}
\textbf{a}, Perfil diurno médio de potência, em pu da capacidade instalada, na
janela comum aos dois ativos. Em azul a geração verificada, em vermelho a
energia cortada.
\textbf{b}, Custo Marginal de Operação do subsistema Sudeste ponderado pela
energia, durante os intervalos de corte e durante a geração.
\textbf{c}, Taxa de corte mensal dos dois ativos.
\textbf{d}, Pixels Sentinel-2 ocupados pela área de módulos, comparados à
instalação mediana de geração distribuída, em escala logarítmica, com área de
módulos estimada a 5,12 m\textsuperscript{2}/kWp.
\textbf{e}, Taxa de corte contra capacidade renovável variável num raio de
100 km, para os conjuntos com apuração completa na janela de referência da
frota, com os dois alvos em destaque.}
\label{fig:caso}
\end{figure*}

A taxa de corte é de 26,8\% contra 8,8\%. A energia cortada equivale a cerca de
R\$ 64 milhões contra pouco mais de R\$ 1 milhão, valorada ao Custo Marginal de
Operação. Nos dois ativos o CMO ponderado durante o corte é menor que durante a
geração, o que é consistente com restrição de razão energética. O
capítulo~\ref{cap:despacho} examina esse ponto no despacho.

\campo{Uma correção de atribuição} O campo \texttt{nom\_\allowbreak agenteproprietario} do
ONS identifica o agente representante, e não o proprietário. Ele não sustenta a
leitura de que Três Marias 3 seja ativo estatal. O SIGA registra sete
proprietários privados sob regime de autoprodução. Não há usina solar estatal no
parque despachado, e o eixo público contra privado não é testável nesta amostra.

\conclusao{Nenhuma das cinco covariáveis locais reproduz a diferença do par na
frota: nem carga do ponto de conexão, nem número de conjuntos, nem densidade
regional de renovável, nem tensão, nem capacidade própria.}

Duas explicações locais eram plausíveis à vista do par: barramento compartilhado
e densidade regional de renovável. Nenhuma sobrevive ao teste na frota. Uma
explicação consistente com um par de ativos não se generaliza a mecanismo sem
esse teste.

\begin{objecao}
Dois ativos não são amostra, e o par foi escolhido por contraste. As cinco
covariáveis foram testadas sem correção de multiplicidade; como nenhuma
resultou significativa, a correção só tornaria a conclusão mais conservadora. O
CMO é proxy de valor sistêmico, e não preço de liquidação nem prejuízo líquido.
\end{objecao}

\campo{O que estava em dado aberto e foi afirmado em falta} A versão anterior
desta análise concluía que a topologia de transmissão não está publicada. Ela
está. O SIGET publica o grafo linha a linha, com 1.166 linhas e 619 subestações
com código ONS, tensão e comprimento, e o pacote \texttt{fator-\allowbreak capacidade-2}
traz a coordenada do ponto de conexão de cada usina. O join cobre 168 dos 229
conjuntos despachados. O limite de transferência também está publicado, no
pacote \texttt{linha-\allowbreak transmissao}, com capacidade operativa em MVA e variantes
sazonais, além de resistência e reatância. Permanece indisponível apenas qual
restrição está ativa em cada intervalo. As cinco covariáveis testadas são locais
ao ponto e não usam a estrutura da rede, de modo que o resultado nulo não se
estende a ela. O capítulo~\ref{cap:rede} a testa.

\section{Cronologia do corte e da energização do entorno}
\label{sec:fig4}

\begin{ficha}
  \item[Janela]      2023 a julho de 2026; entorno desde 2021
  \item[Amostra]     43 e 24 meses, por ativo
  \item[Unidade]     conjunto por mês
  \item[Fonte]       ONS, fator de capacidade e \textit{constrained-off}
  \item[Teste]       nenhum teste formal
\end{ficha}

A apuração de corte começa em abril de 2024, e a onda de construção em Jaíba vai
de 2021 a 2024. A série de corte não cobre o período em que a hipótese de
saturação do entorno se manifestaria. O fator de capacidade existe desde a
entrada em operação e cai quando a geração é limitada:

\[ FC_{\text{verificado}}(t)=\frac{E_{\text{gerada}}(t)}{P_{\text{inst}}\cdot\Delta t} \]

A partir de abril de 2024 a apuração permite recompor o potencial:

\[ FC_{\text{potencial}} = FC_{\text{verificado}} + FC_{\text{cortado}} \]

Se a queda for de restrição de rede, o potencial recomposto permanece estável
enquanto o verificado cai. Se for recurso solar ou degradação, os dois caem
juntos.

\begin{figure*}[t]
\centering
\includegraphics[width=\linewidth]{../figures/fig4_cronologia.pdf}
\caption{\textbf{A queda do fator de capacidade de Sol do Cerrado acompanha a
energização do entorno; o controle não mostra o mesmo.}
\textbf{a}, Jaíba, Minas Gerais. Em área cinza a capacidade fotovoltaica
acumulada em operação num raio de 100 km, no eixo esquerdo; em linha o fator de
capacidade mensal verificado de Sol do Cerrado, no eixo direito; em faixa
vermelha a parcela cortada apurada pelo ONS, disponível a partir de abril de
2024.
\textbf{b}, São Gonçalo do Abaeté, Minas Gerais, com a mesma construção para
Três Marias 3.
\textbf{c}, Fator de capacidade mensal de Sol do Cerrado em 2023 e 2025, mês
contra mês, para separar restrição de sazonalidade.
\textbf{d}, Potência acumulada de geração distribuída nos dois municípios. Os
20,4 MW de Jaíba equivalem a 3,0\% da capacidade de Sol do Cerrado.}
\label{fig:cronologia}
\end{figure*}

\begin{objecao}
A decomposição separa restrição de degradação, e não separa restrição de
variação climática. Se 2025 tiver sido menos ensolarado que 2023, a queda não
informa sobre a rede. O fator de capacidade responde a várias causas ao mesmo
tempo, e a decomposição só vale onde a apuração existe. A comparação é pareada e
sem aleatorização, como na seção~\ref{sec:fig3}.
\end{objecao}

O capítulo~\ref{cap:calibracao} fecha essa lacuna, medindo a irradiância
diretamente sobre cada usina.

\campo{Por que o desenho não estima efeito de tratamento} Cada vizinho energiza
em data distinta, o que caracteriza adoção escalonada. Callaway e Sant'Anna
\cite{callaway2021} mostram que o estimador de diferenças em diferenças com
efeitos fixos duplos fica enviesado nesse caso, por ponderar comparações entre
unidades já tratadas. A leitura é descritiva e pareada contra um único
controle, o que evita o problema e também não estima efeito. Estimá-lo exigiria
o estimador por grupo e período daqueles autores, e uma definição de tratamento
exógena que a seção~\ref{sec:fig2} mostra não existir do lado da geração
distribuída.
